\section{约束子空间}
\label{sec:sobolev-constrained-subspaces}

许多问题要求函数满足边界条件，例如在区间两端取零，或在区域边界上取指定的值。Sobolev函数是几乎处处等价类，而边界是零测集，任意修改边界上的值不改变元素，所以“边界值”不能取一个代表直接读出。本节先用紧支撑光滑函数的闭包定义零边界子空间，它不涉及任何边界值；再在一维用绝对连续代表读出端点值，在高维用迹映射表达边界值。最后研究零边界与零均值约束怎样排除常数方向，使导数单独给出等价范数。

\subsection{紧支撑光滑函数的闭包}
\label{sec:sobolev-compact-closure}

紧支撑光滑函数在边界附近恒为零。若一个Sobolev函数能被这样的函数在Sobolev范数下任意逼近，它就在极限意义上“在边界附近为零”。把这一想法直接取作定义。

\begin{definition}[零边界Sobolev空间]
\label{def:sobolev-w0}
设$\Omega\subset\mathbb R^d$为开集，$k\geq1$为整数，$1\leq p<\infty$。把$C_c^\infty(\Omega)$在$W^{k,p}(\Omega)$范数下的闭包记作$W_0^{k,p}(\Omega)$，称为\term{零边界Sobolev空间}（Sobolev space with zero boundary values），$p=2$时记作$H_0^k(\Omega)$。
\end{definition}

这里的闭包是第~\ref{chap:topological-spaces}章的概念：在由Sobolev范数给出的度量下，每个邻域都与$C_c^\infty(\Omega)$相交的元素全体。度量空间中，它恰好由$C_c^\infty(\Omega)$中序列的极限组成。若$\varphi_n\to f$、$\psi_n\to g$，则对实数$a,b$有$a\varphi_n+b\psi_n\to af+bg$，而$a\varphi_n+b\psi_n$仍是紧支撑光滑函数，所以闭包对线性组合封闭。闭包本身是闭集，于是$W_0^{k,p}(\Omega)$是$W^{k,p}(\Omega)$的闭线性子空间，由定理~\ref{thm:sobolev-complete}，它本身也是Banach空间。

定义只用到逼近，不需要区域边界的任何正则性。两个问题需要分别核对。第一，$W_0^{k,p}(\Omega)$一般是真子空间，下一小节将证明$(0,1)$上的常数$1$不属于$H_0^1(0,1)$。在全空间则没有边界：取在$|x|\leq1$上为$1$、在$|x|\geq2$外为零的截断函数$\chi$，$\chi(x/R)f$在$R\to\infty$时于$W^{k,p}(\mathbb R^d)$中趋于$f$，再对它磨光，可知$W_0^{k,p}(\mathbb R^d)=W^{k,p}(\mathbb R^d)$。第二，定义限于$p<\infty$。在$W^{1,\infty}$范数下，紧支撑光滑函数的极限具有一致连续的导数，第~\ref{sec:sobolev-smooth-density}节末尾的$|x-1/2|$与之距离至少为$1$；所以$p=\infty$时的闭包远小于“零边界的$W^{1,\infty}$函数”，有限指数的结论不能照搬。

\subsection{端点值}
\label{sec:sobolev-endpoint-values}

一维时，边界只有两个端点，而第~\ref{chap:smoothing-weak-derivatives}章已经证明$v\in H^1(0,1)$有绝对连续代表$\tilde v$，并且$\tilde v(y)-\tilde v(x)=\int_x^yv'(t)\,\mathrm dt$。端点值$\tilde v(0)$与$\tilde v(1)$由这个代表确定。我们需要知道它们是否受$H^1$范数控制。

把上式对$x\in(0,1)$积分。由于区间长度为$1$，
\[
\tilde v(y)=\int_0^1\tilde v(x)\,\mathrm dx+\int_0^1\int_x^yv'(t)\,\mathrm dt\,\mathrm dx,\qquad
|\tilde v(y)|\leq\int_0^1|v|\,\mathrm dx+\int_0^1|v'|\,\mathrm dt .
\]
再由Cauchy--Schwarz不等式，$\int_0^1|v|\leq\lVert v\rVert_2$，$\int_0^1|v'|\leq\lVert v'\rVert_2$。于是对每个$y\in[0,1]$，特别是两个端点，
\begin{equation}
|\tilde v(y)|\leq\lVert v\rVert_2+\lVert v'\rVert_2\leq\sqrt2\,\lVert v\rVert_{H^1(0,1)} .
\label{eq:sobolev-endpoint-bound}
\end{equation}
所以映射$v\mapsto\tilde v(0)$是线性的，并且对$H^1$范数连续；$v\mapsto\tilde v(1)$同理。式\eqref{eq:sobolev-endpoint-bound}还说明，$H^1$收敛蕴含连续代表的一致收敛。

只有$L^2$控制是不够的。$v_n(x)=\max(1-nx,0)$满足$v_n(0)=1$，而$\lVert v_n\rVert_2^2=\int_0^{1/n}(1-nx)^2\,\mathrm dx=1/(3n)\to0$。端点值能被控制，靠的是导数项：$\lVert v_n'\rVert_2=\sqrt n$随$n$增大。

现在把估计用于闭包。每个$\varphi\in C_c^\infty(0,1)$在端点为零。若$\varphi_n\to v$在$H^1$中成立，由式\eqref{eq:sobolev-endpoint-bound}，$\tilde v(0)=\lim\varphi_n(0)=0$，同理$\tilde v(1)=0$。所以$H_0^1(0,1)=W_0^{1,2}(0,1)$中的函数都有零端点值。常数$1$不满足这一条件，并且对每个$\varphi\in C_c^\infty(0,1)$，
\[
1=|1-\varphi(0)|\leq\sqrt2\,\lVert1-\varphi\rVert_{H^1},
\]
它到紧支撑光滑函数的$H^1$距离至少是$1/\sqrt2$。反过来，连续代表在两端为零的$H^1(0,1)$函数都属于$H_0^1(0,1)$，这是下一小节零迹刻画在一维的情形。

\subsection{迹映射}
\label{sec:sobolev-trace}

高维时，$W^{1,p}$函数一般没有连续代表，第~\ref{sec:sobolev-continuous-embeddings}节将给出在边界点附近无界的例子，因此不能逐点读取边界值。我们改为先对光滑到边界的函数直接限制到边界，证明限制的边界$L^p$大小受Sobolev范数控制，再用逼近推广。

\paragraph{半空间上的估计}
设$f\in C_c^1(\mathbb R^d)$，考虑它在闭半空间$\overline{\mathbb R^d_+}$上的部分。固定$x'$，沿$x_d$方向从边界点$(x',0)$向内积分。$t$足够大时$f(x',t)=0$，所以$p=1$时
\[
|f(x',0)|=\Bigl|\int_0^\infty\partial_df(x',t)\,\mathrm dt\Bigr|\leq\int_0^\infty|\partial_df(x',t)|\,\mathrm dt .
\]
$p>1$时$|f|^p$连续可微，导数的绝对值不超过$p|f|^{p-1}|\partial_df|$。由Young不等式$ab\leq a^{p'}/p'+b^p/p$（取$a=|f|^{p-1}$，$b=|\partial_df|$，$p'=p/(p-1)$），
\[
|f(x',0)|^p=-\int_0^\infty\partial_t\bigl(|f(x',t)|^p\bigr)\,\mathrm dt\leq\int_0^\infty\bigl((p-1)|f|^p+|\partial_df|^p\bigr)\,\mathrm dt .
\]
两种情形都对$x'\in\mathbb R^{d-1}$积分，得到
\begin{equation}
\int_{\mathbb R^{d-1}}|f(x',0)|^p\,\mathrm dx'\leq\max(p-1,1)\,\lVert f\rVert_{W^{1,p}(\mathbb R^d_+)}^p .
\label{eq:sobolev-half-space-trace}
\end{equation}
边界值的$L^p$大小被体积上函数与法向导数的大小控制。这与端点估计是同一机制：沿垂直于边界的每条直线，问题都是一维的。

\paragraph{Lipschitz区域上的迹}
在有界Lipschitz区域上，边界需要配上\term{表面测度}（surface measure）$\mathrm dS$。在定义~\ref{def:sobolev-lipschitz-domain}的局部图像$x_d=\gamma(x')$上，取$\mathrm dS=\sqrt{1+|\nabla\gamma(x')|^2}\,\mathrm dx'$。Lipschitz函数几乎处处可微且梯度有界，所以这个因子介于$1$与$\sqrt{1+L_\gamma^2}$之间，$L_\gamma$为$\gamma$的Lipschitz常数；不同局部图像在重叠处给出同一测度，这一点这里不加证明地使用。记$L^p(\partial\Omega)=L^p(\partial\Omega,\mathrm dS)$。

\begin{theorem}[迹定理]
\label{thm:sobolev-trace}
设$\Omega\subset\mathbb R^d$为有界Lipschitz区域，$1\leq p<\infty$。存在常数$C$，使所有$f\in C^\infty(\overline\Omega)$满足
\begin{equation}
\lVert f|_{\partial\Omega}\rVert_{L^p(\partial\Omega)}\leq C\lVert f\rVert_{W^{1,p}(\Omega)} .
\label{eq:sobolev-trace-inequality}
\end{equation}
并且存在唯一的有界线性算子$\operatorname{Tr}:W^{1,p}(\Omega)\to L^p(\partial\Omega)$，使$\operatorname{Tr}f=f|_{\partial\Omega}$对所有$f\in C^\infty(\overline\Omega)$成立。
\end{theorem}

不等式\eqref{eq:sobolev-trace-inequality}的证明思路与延拓定理相同：用有限个边界球和一个内部开集覆盖$\overline\Omega$，取从属的单位分解把$f$分成有限块；内部那块在边界上为零；每个边界块经拉直变成半空间上的紧支撑函数，套用式\eqref{eq:sobolev-half-space-trace}，再把拉直变换与表面测度因子带来的有界倍数计入常数。这里不给出细节。

\begin{proof}[算子的构造]
设不等式\eqref{eq:sobolev-trace-inequality}成立。对$f\in W^{1,p}(\Omega)$，由定理~\ref{thm:sobolev-smooth-up-to-boundary}取$f_n\in C^\infty(\overline\Omega)$在$W^{1,p}$中趋于$f$。由式\eqref{eq:sobolev-trace-inequality}，
\[
\lVert f_n|_{\partial\Omega}-f_m|_{\partial\Omega}\rVert_{L^p(\partial\Omega)}\leq C\lVert f_n-f_m\rVert_{W^{1,p}(\Omega)},
\]
所以边界限制在$L^p(\partial\Omega)$中构成Cauchy列，收敛到某个$g$。若$h_n$是另一个逼近列，交错排列的$(f_1,h_1,f_2,h_2,\ldots)$仍趋于$f$，其边界限制也收敛，于是两个子列的极限相同，$g$不依赖逼近列。令$\operatorname{Tr}f=g$。线性由极限的线性得到，$\lVert\operatorname{Tr}f\rVert_{L^p(\partial\Omega)}\leq C\lVert f\rVert_{W^{1,p}}$由不等式取极限得到。两个有界线性算子若在稠密子集上相同，便处处相同，这给出唯一性。
\end{proof}

称$\operatorname{Tr}f$为$f$的\term{边界迹}（boundary trace）。它是$L^p(\partial\Omega)$中的等价类，而不是任选一个代表后逐点读出的边界值：修改$f$在零测集$\partial\Omega$上的值不改变$f$，也不改变$\operatorname{Tr}f$；修改代表在区域内部某个零测集上的值同样不影响它。迹由“函数在靠近边界处的平均行为”决定，这正是逼近构造所表达的。

零迹与零边界子空间在一阶情形可以对应起来。$C_c^\infty(\Omega)$中的函数在边界附近为零，迹为零；$\operatorname{Tr}$连续，零迹在$W^{1,p}$极限下保持，所以$W_0^{1,p}(\Omega)\subset\ker\operatorname{Tr}$。反向包含在有界Lipschitz区域上也成立，因而
\[
W_0^{1,p}(\Omega)=\{f\in W^{1,p}(\Omega):\operatorname{Tr}f=0\}.
\]
反向包含的证明思路是：在每个局部图像中把零迹函数沿远离边界的方向推进一小段，乘以在边界附近为零的截断函数，再磨光；零迹保证截断造成的误差随推进距离趋于零。这里不给出细节。

这一对应不能直接推广为高阶零迹的完整刻画。在$(0,1)$上取$f(x)=x(1-x)\in H^2(0,1)$，它的端点值为零，但$f'(0)=1$。对$\varphi\in C_c^\infty(0,1)$，把式\eqref{eq:sobolev-endpoint-bound}用于$(f-\varphi)'$，得$1=|f'(0)-\varphi'(0)|\leq\sqrt2\lVert f-\varphi\rVert_{H^2}$，所以$f\notin H_0^2(0,1)$。高阶闭包还要求法向导数的边界值为零，刻画它需要对导数的迹另作讨论。

\subsection{约束子空间中的导数范数}
\label{sec:sobolev-derivative-norms}

第~\ref{sec:sobolev-oscillation}小节指出，常数的导数为零，因此在整个$W^{1,p}(\Omega)$上，导数的大小不能控制函数本身。记$|\nabla f|$为弱梯度$(\partial_1f,\ldots,\partial_df)$的欧氏长度，$\lVert\nabla f\rVert_p$为它的$L^p$范数。逐点有$|\partial_if|\leq|\nabla f|\leq\sum_i|\partial_if|$，所以$\lVert\nabla f\rVert_p$与$(\sum_i\lVert\partial_if\rVert_p^p)^{1/p}$只相差依赖$d,p$的倍数。我们要回答：在排除常数方向的子空间上，是否存在常数$C$使$\lVert f\rVert_p\leq C\lVert\nabla f\rVert_p$。零边界与零均值是两种排除常数的方式。

\paragraph{零边界情形}
形如$\lVert f\rVert_p\leq C\lVert\nabla f\rVert_p$的估计称为\term{Poincaré不等式}（Poincaré inequality）。零边界条件下，它不需要区域连通或边界正则，只需要区域在一个方向上有界。

\begin{theorem}[零边界的Poincaré不等式]
\label{thm:sobolev-poincare-zero}
设$\Omega\subset\mathbb R^d$为开集，含于带状区域$\{x:a<x_1<a+L\}$，$1\leq p<\infty$。则对所有$f\in W_0^{1,p}(\Omega)$，
\[
\lVert f\rVert_p\leq\frac{L}{p^{1/p}}\lVert\partial_1f\rVert_p\leq\frac{L}{p^{1/p}}\lVert\nabla f\rVert_p .
\]
\end{theorem}

\begin{proof}
先设$f\in C_c^\infty(\Omega)$，并把它在$\Omega$外补成零。记$x=(x_1,x'')$。由于$f(a,x'')=0$，
\begin{equation*}
\begin{aligned}
|f(x_1,x'')|&=\Bigl|\int_a^{x_1}\partial_1f(t,x'')\,\mathrm dt\Bigr|\\
&\leq(x_1-a)^{1-1/p}\Bigl(\int_a^{a+L}|\partial_1f(t,x'')|^p\,\mathrm dt\Bigr)^{1/p},
\end{aligned}
\end{equation*}
第二步是Hölder不等式。取$p$次方，对$x_1\in(a,a+L)$积分，$\int_a^{a+L}(x_1-a)^{p-1}\,\mathrm dx_1=L^p/p$；再对$x''$积分，得$\lVert f\rVert_p^p\leq(L^p/p)\lVert\partial_1f\rVert_p^p$。两边都对$W^{1,p}$范数连续，由$W_0^{1,p}(\Omega)$的定义取极限即得一般情形。
\end{proof}

\paragraph{零均值情形}
零均值$f_\Omega=|\Omega|^{-1}\int_\Omega f\,\mathrm dx=0$同样排除了非零常数，但它是整体条件，需要区域连通，否则可以在不同分支上取不同常数。我们先在凸区域上用积分估计完整证明，再说明一般区域怎样局部拼接。下面的引理用一点的值与区域平均之差来读取梯度的作用。

\begin{lemma}[凸区域上的平均估计]
\label{lem:sobolev-convex-average}
设$W\subset\mathbb R^d$为有界凸开集，直径为$D$，$f$在$\overline W$的一个邻域上连续可微。则对每个$x\in\overline W$，
\begin{equation}
|f(x)-f_W|\leq\frac{D^d}{d\,|W|}\int_W\frac{|\nabla f(z)|}{|x-z|^{d-1}}\,\mathrm dz .
\label{eq:sobolev-convex-average}
\end{equation}
\end{lemma}

\begin{proof}
对$y\in W$，记$\omega=(y-x)/|y-x|$。凸性保证除端点$x$外，线段$[x,y]$都落在$W$中，所以$|f(y)-f(x)|\leq\int_0^{|y-x|}G(x+t\omega)\,\mathrm dt$，其中$G=|\nabla f|\mathbf 1_W$。由$f(x)-f_W=|W|^{-1}\int_W(f(x)-f(y))\,\mathrm dy$，
\[
|f(x)-f_W|\leq\frac1{|W|}\int_W\int_0^\infty G(x+t\omega)\,\mathrm dt\,\mathrm dy .
\]
以$x$为中心用球坐标$y=x+\rho\omega$，$\mathrm dy=\rho^{d-1}\mathrm d\rho\,\mathrm d\sigma(\omega)$，$\rho\leq D$。内层积分只依赖方向$\omega$，对$\rho$积分给出因子$D^d/d$。剩下的$\int_{S^{d-1}}\int_0^\infty G(x+t\omega)\,\mathrm dt\,\mathrm d\sigma(\omega)$再按球坐标还原为$\int G(z)|z-x|^{1-d}\,\mathrm dz$，即得式\eqref{eq:sobolev-convex-average}。
\end{proof}

\begin{theorem}[凸区域上的零均值Poincaré不等式]
\label{thm:sobolev-poincare-convex}
设$W\subset\mathbb R^d$为有界凸开集，直径为$D$，$1\leq p<\infty$。则对所有$f\in W^{1,p}(W)$，
\[
\lVert f-f_W\rVert_p\leq\frac{|B(0,D)|}{|W|}\,D\,\lVert\nabla f\rVert_p .
\]
\end{theorem}

\begin{proof}
先设$f\in C^\infty(\overline W)$。记$Tg(x)=\int_W g(z)|x-z|^{1-d}\,\mathrm dz$。对每个$x\in W$，$W\subset B(x,D)$，所以$\int_W|x-z|^{1-d}\,\mathrm dz\leq\int_{B(x,D)}|x-z|^{1-d}\,\mathrm dz=|S^{d-1}|D$，其中$|S^{d-1}|$为单位球面的面积；对$z$固定、$x$变化也一样。把核写成$|x-z|^{(1-d)/p'}\cdot|x-z|^{(1-d)/p}$，用Hölder不等式，
\[
|Tg(x)|^p\leq\bigl(|S^{d-1}|D\bigr)^{p-1}\int_W|g(z)|^p|x-z|^{1-d}\,\mathrm dz ,
\]
再对$x$积分并交换积分次序，得$\lVert Tg\rVert_p\leq|S^{d-1}|D\,\lVert g\rVert_p$。把它与式\eqref{eq:sobolev-convex-average}合并，并用$|B(0,D)|=|S^{d-1}|D^d/d$，即得所需估计。有界凸开集是Lipschitz区域，由定理~\ref{thm:sobolev-smooth-up-to-boundary}，$C^\infty(\overline W)$在$W^{1,p}(W)$中稠密；均值满足$|f_W-g_W|\leq|W|^{-1/p}\lVert f-g\rVert_p$，两边都对$W^{1,p}$范数连续，取极限得一般情形。
\end{proof}

对一般的有界连通Lipschitz区域$\Omega$，零均值的Poincaré不等式$\lVert f-f_\Omega\rVert_p\leq C\lVert\nabla f\rVert_p$同样成立，$1\leq p<\infty$，常数依赖$\Omega$与$p$。局部拼接的思路如下：用有限个已有局部估计的小区域覆盖$\Omega$，内部用球，边界附近用拉直后的柱形区域；连通性保证任意两块能由一串依次重叠的小区域相连，重叠部分测度为正，于是相邻两块均值之差被梯度控制，逐块传递得到整体估计。这里不给出完整证明。第~\ref{sec:sobolev-compact-lq}小节的紧嵌入还给出另一条路线：若估计不成立，可取$\lVert f_n\rVert_p=1$、均值为零、$\lVert\nabla f_n\rVert_p\to0$的序列，紧嵌入给出$L^p$收敛子列，其极限梯度为零，在连通区域上只能是常数，零均值迫使它为零，与范数为$1$矛盾。

\paragraph{条件为什么不能少}
连通性不能去掉：在$(0,1)\cup(2,3)$上，第一段取$-1$、第二段取$1$的函数均值为零、弱导数为零，却不是零。约束不能去掉：常数函数本身就是反例。零边界情形中区域有界的方向也不能去掉：在$\mathbb R$上取非零$\psi\in C_c^\infty(\mathbb R)$，令$\psi_R(x)=\psi(x/R)$，则$\lVert\psi_R\rVert_2=R^{1/2}\lVert\psi\rVert_2$，$\lVert\psi_R'\rVert_2=R^{-1/2}\lVert\psi'\rVert_2$，比值随$R$无界。

\paragraph{导数给出的等价范数}
Poincaré不等式把约束和导数范数联系起来。在定理~\ref{thm:sobolev-poincare-zero}的条件下，对$f\in W_0^{1,p}(\Omega)$，
\[
\frac{1}{d^{1-1/p}}\lVert\nabla f\rVert_p\leq\lVert f\rVert_{W^{1,p}}\leq\Bigl(\frac{L^p}{p}+d\Bigr)^{1/p}\lVert\nabla f\rVert_p .
\]
左边来自$\lVert\nabla f\rVert_p\leq\sum_i\lVert\partial_if\rVert_p$与$\ell^p$、$\ell^1$范数的比较，右边来自Poincaré不等式与$|\partial_if|\leq|\nabla f|$。于是在零边界子空间上，$\lVert\nabla f\rVert_p$就是一个等价范数；$p=2$时，$\int_\Omega\nabla f\cdot\nabla g$是$H_0^1(\Omega)$上与原内积等价的内积。在连通Lipschitz区域的零均值子空间上，结论相同。在整个$W^{1,p}(\Omega)$上则不行，二者的差别恰好是常数方向。
